% BHU PYQ -- Transcribed & Corrected
% Paper: ENGAE-21 / ENGAE21: Listening & Speaking Skills
% Exam: Under Graduate Semester II Examination, 2024-25 (NEP)
% Category: Ability Enhancement Course (AEC)
% Department: English

\documentclass[11pt,a4paper]{article}
\usepackage[a4paper,top=2.5cm,bottom=2.5cm,left=2.8cm,right=2.8cm,headheight=14pt]{geometry}
\usepackage{amsmath,amssymb}
\usepackage[T1]{fontenc}
\usepackage[utf8]{inputenc}
\usepackage{lmodern,microtype,fancyhdr,enumitem,hyperref,lastpage,parskip}

\hypersetup{
    colorlinks=true,
    urlcolor=black,
    linkcolor=black,
    pdftitle={ENGAE21: Listening & Speaking Skills -- BHU Sem II 2024-25 (NEP AEC)}
}

\pagestyle{fancy}
\fancyhf{}
\renewcommand{\headrulewidth}{0.4pt}
\renewcommand{\footrulewidth}{0.4pt}
\fancyhead[L]{\small   \textit{Banaras Hindu University}}
\fancyhead[R]{\small   \textit{ENGAE-21: Listening \& Speaking Skills (AEC)}}
\fancyfoot[C]{\small Page  \thepage\ of \pageref{LastPage}}


\newlist{parts}{enumerate}{2}
\setlist[parts,1]{label=(\Alph*),leftmargin=*,itemsep=3pt,topsep=2pt}

\newcommand{\pts}[1]{\hfill   \textit{[#1]}}


\begin{document}


\begin{center}
  {\large\bfseries BANARAS HINDU UNIVERSITY}\\[2pt]
  {
\normalsize Under Graduate --- Semester II Examination, 2024-25 (NEP)}\\[6pt]
  \rule{0.8\linewidth}{0.5pt}\\[6pt]
  {\large\bfseries Ability Enhancement Course (AEC)}\\[4pt]
  {\large\bfseries Paper Code: ENGAE-21 : Listening \& Speaking Skills}\\[4pt]
  {
\normalsize Time: 1.5 Hours\hfill Maximum Marks: 70}\\[2pt]
  {\small REF NO. CE/Conf./D-II/74}
\end{center}

\rule{\linewidth}{0.4pt}

\smallskip



\noindent   \textit{Instructions: Answer all 70 Multiple Choice Questions. Each question carries equal marks (1 mark each).}

\bigskip


\noindent   \textbf{Question 1.} What is the main goal of group discussions in listening exercises?\pts{1}

\begin{parts}
  \item To explore varied views through shared listening
  \item To test writing skills using discussion topics
  \item To assess pronunciation and fluency
  \item To promote silence and observation
\end{parts}

\medskip


\noindent   \textbf{Question 2.} Reflective writing on listening helps students:\pts{1}

\begin{parts}
  \item Improve punctuation and formatting
  \item Practice memorising quotes
  \item Understand how listening affects outcomes
  \item Take longer notes during lectures
\end{parts}

\medskip


\noindent   \textbf{Question 3.} Active listening primarily involves:\pts{1}

\begin{parts}
  \item Glancing while multitasking
  \item Asking frequent off-topic questions
  \item Paying full attention and seeking to understand
  \item Memorising every spoken word
\end{parts}

\medskip


\noindent   \textbf{Question 4.} Listening to diverse accents helps students to:\pts{1}

\begin{parts}
  \item Adopt a new accent
  \item Improve comprehension of varied speech styles
  \item Write faster transcriptions
  \item Avoid listening to local speakers
\end{parts}

\medskip


\noindent   \textbf{Question 5.} In classroom listening, internal distractions may include:\pts{1}

\begin{parts}
  \item Noise from outside
  \item Loud air conditioner
  \item Personal thoughts and anxieties
  \item Visual slides
\end{parts}

\medskip


\noindent   \textbf{Question 6.} What can interfere with effective listening?\pts{1}

\begin{parts}
  \item Calm surroundings and interest
  \item Summarising during pauses
  \item Respect for differing opinions
  \item Emotional distractions and noise
\end{parts}

\medskip


\noindent   \textbf{Question 7.} In classroom listening, external noise may include:\pts{1}

\begin{parts}
  \item A student's inner thoughts
  \item Nearby construction sounds
  \item Poor handwriting on the board
  \item Loud font usage in slides
\end{parts}

\medskip


\noindent   \textbf{Question 8.} Which task requires tracking listening habits?\pts{1}

\begin{parts}
  \item Grammar journal keeping
  \item Maintaining a listening logbook
  \item Reading newspaper columns
  \item Daily schedule planning
\end{parts}

\medskip


\noindent   \textbf{Question 9.} A listening journal is meant to:\pts{1}

\begin{parts}
  \item Record thoughts and reflect on listening experiences
  \item Practice cursive handwriting
  \item Store lecture outlines
  \item Track exam marks
\end{parts}

\medskip


\noindent   \textbf{Question 10.} Daily reflection on listening helps to:\pts{1}

\begin{parts}
  \item Increase writing speed
  \item Develop self-awareness as a listener
  \item Eliminate the need for note-taking
  \item Memorize textbook facts
\end{parts}

\medskip


\noindent   \textbf{Question 11.} One strategy to improve listening in noisy places is:\pts{1}

\begin{parts}
  \item Tune out the speaker completely
  \item Focus only on background music
  \item Rely on facial cues and summary skills
  \item Whisper responses back to yourself
\end{parts}

\medskip


\noindent   \textbf{Question 12.} Which action is not part of active listening?\pts{1}

\begin{parts}
  \item Restating the main point
  \item Asking follow-up questions
  \item Interrupting to share your opinion
  \item Showing interest through body language
\end{parts}

\medskip


\noindent   \textbf{Question 13.} Exploring mediums like film and radio is useful because:\pts{1}

\begin{parts}
  \item They simplify complex ideas
  \item They expose students to varied speech styles
  \item They require less attention
  \item They replace reading books
\end{parts}

\medskip


\noindent   \textbf{Question 14.} What skill is built through story-sharing with peers?\pts{1}

\begin{parts}
  \item Empathetic listening with attention
  \item Memorisation for exams
  \item Story rewriting techniques
  \item Visualisation for recall
\end{parts}

\medskip


\noindent   \textbf{Question 15.} A key challenge in real-time listening is:\pts{1}

\begin{parts}
  \item Limited vocabulary
  \item Keeping pace with spoken words
  \item Too many written questions
  \item Short video length
\end{parts}

\medskip


\noindent   \textbf{Question 16.} In the Medium Exploration task, students:\pts{1}

\begin{parts}
  \item Write essays about TV actors
  \item Record phone calls for grading
  \item Present insights using creative media
  \item Take notes on textbook margins
\end{parts}

\medskip


\noindent   \textbf{Question 17.} Why use multiple formats (audio, video, live speeches) in listening practice?\pts{1}

\begin{parts}
  \item To confuse the listener
  \item To cater to diverse learning preferences and styles
  \item To reduce overall practice time
  \item To avoid using textbooks
\end{parts}

\medskip


\noindent   \textbf{Question 18.} In an Empathy Circle, students are expected to practice:\pts{1}

\begin{parts}
  \item Active listening and empathetic responses
  \item Quick thinking under pressure
  \item Formal debate
  \item Nonverbal imitation
\end{parts}

\medskip


\noindent   \textbf{Question 19.} What is the Confidence Circle activity primarily designed to improve?\pts{1}

\begin{parts}
  \item Vocabulary retention
  \item Argument development
  \item Confidence in speaking and peer support
  \item Visual aids for speaking
\end{parts}

\medskip


\noindent   \textbf{Question 20.} What is the goal of the Visualization Exercise in public speaking training?\pts{1}

\begin{parts}
  \item Prepare slides for presentation
  \item Mentally rehearse a confident and successful speech
  \item Create a speech outline
  \item Evaluate body language
\end{parts}

\medskip


\noindent   \textbf{Question 21.} Which of the following best describes the purpose of Listening Drills?\pts{1}

\begin{parts}
  \item Expanding vocabulary through repetition
  \item Practicing pronunciation
  \item Memorizing content
  \item Enhancing summarization and inference skills
\end{parts}

\medskip


\noindent   \textbf{Question 22.} What skill is primarily targeted during Gestural Analysis?\pts{1}

\begin{parts}
  \item Sentence structuring
  \item Understanding nonverbal communication cues
  \item Vocal projection
  \item Brainstorming techniques
\end{parts}

\medskip


\noindent   \textbf{Question 23.} Which classroom activity involves interpreting and providing nonverbal cues through acting?\pts{1}

\begin{parts}
  \item Storytelling Workshop
  \item Confidence Circle
  \item Gestural Feedback Role-Play
  \item Listening Drill
\end{parts}

\medskip


\noindent   \textbf{Question 24.} Why is gestural feedback important in public speaking?\pts{1}

\begin{parts}
  \item It helps reduce speech time
  \item It improves clarity and strengthens communication
  \item It enhances memorization
  \item It replaces verbal content entirely
\end{parts}

\medskip


\noindent   \textbf{Question 25.} What does the Vocal Warm-Up Routine primarily improve?\pts{1}

\begin{parts}
  \item Vocabulary building
  \item Breath control, resonance and articulation
  \item Topic generation
  \item Note-taking
\end{parts}

\medskip


\noindent   \textbf{Question 26.} Which of the following is a key focus in speech practice sessions?\pts{1}

\begin{parts}
  \item Facial expressions
  \item Background music selection
  \item Slide organization
  \item Volume, tone and clarity
\end{parts}

\medskip


\noindent   \textbf{Question 27.} What is the objective of the Storytelling Workshop in this course?\pts{1}

\begin{parts}
  \item Improve handwriting and fluency
  \item Study mythological themes
  \item Learn to connect with the audience
  \item Practice stage direction
\end{parts}

\medskip


\noindent   \textbf{Question 28.} How do interactive speeches aim to engage the audience?\pts{1}

\begin{parts}
  \item By avoiding direct interaction
  \item Through questions, polls or invitations to respond
  \item By using only visual aids
  \item Through scripted monologues and performances
\end{parts}

\medskip


\noindent   \textbf{Question 29.} Which activity directly helps students overcome stage fright?\pts{1}

\begin{parts}
  \item Listening drills
  \item Story analysis
  \item Confidence Circle
  \item Grammar worksheets
\end{parts}

\medskip


\noindent   \textbf{Question 30.} What is a major benefit of using personal anecdotes in a speech?\pts{1}

\begin{parts}
  \item Ensures logical structure
  \item Broadens vocabulary usage
  \item Relatability with the audience
  \item Lengthens the speech
\end{parts}

\medskip


\noindent   \textbf{Question 31.} Which exercise is designed to promote group bonding and encourage self-expression?\pts{1}

\begin{parts}
  \item Icebreaker
  \item Role-play feedback
  \item Report writing
  \item Pronunciation test
\end{parts}

\medskip


\noindent   \textbf{Question 32.} Gestural feedback can include:\pts{1}

\begin{parts}
  \item Loud vocal emphasis
  \item Repeating key ideas
  \item Facial expressions
  \item Taking notes silently
\end{parts}

\medskip


\noindent   \textbf{Question 33.} Active listening includes which of the following?\pts{1}

\begin{parts}
  \item Ignoring body language
  \item Interrupting with opinions
  \item Memorizing facts and blurting them during the lecture
  \item Summarizing, empathizing and observing responses
\end{parts}

\medskip


\noindent   \textbf{Question 34.} What is the benefit of vocal warm-ups before delivering a speech?\pts{1}

\begin{parts}
  \item Avoids using a microphone
  \item Prepares vocal range and projection
  \item Increases speech duration
  \item Removes the need for body gestures
\end{parts}

\medskip


\noindent   \textbf{Question 35.} What is a likely outcome of consistently practising speech delivery?\pts{1}

\begin{parts}
  \item Improved vocal control and confidence
  \item Decreased spontaneity
  \item Shorter attention spans
  \item Increased filler word use
\end{parts}

\medskip


\noindent   \textbf{Question 36.} Interactive elements in speeches are useful because they:\pts{1}

\begin{parts}
  \item Allow the speaker to pause
  \item Focus only on speaker's views
  \item Eliminate the need for visuals
  \item Boost audience participation and attentiveness
\end{parts}

\medskip


\noindent   \textbf{Question 37.} Which of the following is an active articulator?\pts{1}

\begin{parts}
  \item Teeth
  \item Hard palate
  \item Tongue
  \item Alveolar ridge
\end{parts}

\medskip


\noindent   \textbf{Question 38.} The lungs are responsible for:\pts{1}

\begin{parts}
  \item Articulation
  \item Voice modulation
  \item Providing the air stream
  \item Resonance
\end{parts}

\medskip


\noindent   \textbf{Question 39.} The vocal cords are located in the:\pts{1}

\begin{parts}
  \item Pharynx
  \item Nasal cavity
  \item Larynx
  \item Trachea
\end{parts}

\medskip


\noindent   \textbf{Question 40.} Which of the following is a passive articulator?\pts{1}

\begin{parts}
  \item Lower lip
  \item Soft palate
  \item Uvula
  \item Upper teeth
\end{parts}

\medskip


\noindent   \textbf{Question 41.} Which organ determines the quality of nasal sounds?\pts{1}

\begin{parts}
  \item Tongue
  \item Uvula
  \item Epiglottis
  \item Vocal cords
\end{parts}

\medskip


\noindent   \textbf{Question 42.} Which of these organs is not directly involved in speech production?\pts{1}

\begin{parts}
  \item Tongue
  \item Lungs
  \item Liver
  \item Larynx
\end{parts}

\medskip


\noindent   \textbf{Question 43.} Which pair of speech organs are used to produce bilabial sounds?\pts{1}

\begin{parts}
  \item Tongue and teeth
  \item Upper and lower lips
  \item Tongue and alveolar ridge
  \item Uvula and soft palate
\end{parts}

\medskip


\noindent   \textbf{Question 44.} The roof of the mouth is made up of the:\pts{1}

\begin{parts}
  \item Alveolar ridge
  \item Soft palate and hard palate
  \item Uvula
  \item Nasal cavity
\end{parts}

\medskip


\noindent   \textbf{Question 45.} The glottis is:\pts{1}

\begin{parts}
  \item A speech sound
  \item The space between the vocal cords
  \item Another name for the larynx
  \item A part of the nasal cavity
\end{parts}

\medskip


\noindent   \textbf{Question 46.} Which of the following organs help in voicing (producing voiced sounds)?\pts{1}

\begin{parts}
  \item Nasal cavity
  \item Vocal cords
  \item Tongue
  \item Teeth
\end{parts}

\medskip


\noindent   \textbf{Question 47.} Phonetics is the study of:\pts{1}

\begin{parts}
  \item Word meanings
  \item Sound patterns in a language
  \item Sentence structures
  \item Social use of language
\end{parts}

\medskip


\noindent   \textbf{Question 48.} Which of the following branches of phonetics deals with the physical properties of speech sounds?\pts{1}

\begin{parts}
  \item Auditory phonetics
  \item Acoustic phonetics
  \item Articulatory phonetics
  \item Visual phonetics
\end{parts}

\medskip


\noindent   \textbf{Question 49.} The International Phonetic Alphabet (IPA) is used to:\pts{1}

\begin{parts}
  \item Teach grammar
  \item Represent speech sounds universally
  \item Create poetry
  \item Translate languages
\end{parts}

\medskip


\noindent   \textbf{Question 50.} How many vowel sounds are there in Received Pronunciation (RP) English?\pts{1}

\begin{parts}
  \item 5
  \item 10
  \item 20
  \item 26
\end{parts}

\medskip


\noindent   \textbf{Question 51.} Which of the following sounds is a voiceless consonant?\pts{1}

\begin{parts}
  \item /b/
  \item /d/
  \item /s/
  \item /g/
\end{parts}

\medskip


\noindent   \textbf{Question 52.} Articulatory phonetics is concerned with:\pts{1}

\begin{parts}
  \item How sounds are heard
  \item How sounds are produced by speech organs
  \item The history of sounds
  \item Sound change over time
\end{parts}

\medskip


\noindent   \textbf{Question 53.} Which of the following is a diphthong?\pts{1}

\begin{parts}
  \item /i:/
  \item /e/
  \item /aɪ/
  \item /u/
\end{parts}

\medskip


\noindent   \textbf{Question 54.} Which sound is a voiced dental fricative?\pts{1}

\begin{parts}
  \item /θ/
  \item /ʃ/
  \item /ð/
  \item /f/
\end{parts}

\medskip


\noindent   \textbf{Question 55.} The IPA symbol /ʃ/ corresponds to the sound in:\pts{1}

\begin{parts}
  \item "Go"
  \item "She"
  \item "This"
  \item "Cat"
\end{parts}

\medskip


\noindent   \textbf{Question 56.} Which of the following is not a vowel sound?\pts{1}

\begin{parts}
  \item /eɪ/
  \item /i:/
  \item /z/
  \item /ɔ:/
\end{parts}

\medskip


\noindent   \textbf{Question 57.} What is the primary purpose of visual aids in a presentation?\pts{1}

\begin{parts}
  \item To entertain the audience
  \item To make the speaker look professional
  \item To support and enhance understanding of the message
  \item To reduce the speaker's preparation time
\end{parts}

\medskip


\noindent   \textbf{Question 58.} Which of the following is an example of a presentation tool?\pts{1}

\begin{parts}
  \item Chalkboard
  \item PowerPoint
  \item Poster
  \item All of the above
\end{parts}

\medskip


\noindent   \textbf{Question 59.} Which visual aid is most effective for displaying numerical data clearly?\pts{1}

\begin{parts}
  \item Chart or graph
  \item Video clip
  \item Hand-drawn sketch
  \item Whiteboard
\end{parts}

\medskip


\noindent   \textbf{Question 60.} When using visual aids a speaker should:\pts{1}

\begin{parts}
  \item Read everything written on the slides
  \item Use lots of bright colors and text
  \item Keep text minimal and visuals clear
  \item Avoid rehearsing with the aid
\end{parts}

\medskip


\noindent   \textbf{Question 61.} A good PowerPoint slide typically includes:\pts{1}

\begin{parts}
  \item Paragraphs of text
  \item One main idea and key points
  \item No visuals
  \item Full script of the speech
\end{parts}

\medskip


\noindent   \textbf{Question 62.} Adapting to different speaking contexts means:\pts{1}

\begin{parts}
  \item Using the same speech everywhere
  \item Changing your accent to match the audience
  \item Modifying your tone, language and examples based on the audience and occasion
  \item Speaking faster to save time
\end{parts}

\medskip


\noindent   \textbf{Question 63.} Which factor is most important when addressing a formal audience?\pts{1}

\begin{parts}
  \item Wearing casual clothes
  \item Using slang and humor freely
  \item Maintaining respectful tone and structured delivery
  \item Avoiding eye contact
\end{parts}

\medskip


\noindent   \textbf{Question 64.} How should a speaker adapt when presenting to young children?\pts{1}

\begin{parts}
  \item Use complex vocabulary
  \item Speak in a monotone
  \item Include stories, visuals and simple language
  \item Use technical jargon
\end{parts}

\medskip


\noindent   \textbf{Question 65.} In an online presentation, which of the following is essential?\pts{1}

\begin{parts}
  \item Background music
  \item Strong internet connection and clear visuals
  \item Fast talking
  \item Distracting virtual backgrounds
\end{parts}

\medskip


\noindent   \textbf{Question 66.} What is the role of rehearsal when using visual aids?\pts{1}

\begin{parts}
  \item To memorize the slides
  \item To practice transitions and timing
  \item To avoid looking at the audience
  \item To increase nervousness
\end{parts}

\medskip


\noindent   \textbf{Question 67.} Which element is crucial for effective voice projection in public speaking?\pts{1}

\begin{parts}
  \item Speaking very fast
  \item Diaphragmatic breathing and posture
  \item Shouting
  \item Whispering
\end{parts}

\medskip


\noindent   \textbf{Question 68.} Nonverbal communication includes which of the following elements?\pts{1}

\begin{parts}
  \item Word choice
  \item Grammar and syntax
  \item Eye contact, posture, and facial expressions
  \item Font size on slides
\end{parts}

\medskip


\noindent   \textbf{Question 69.} Which strategy helps in managing public speaking anxiety?\pts{1}

\begin{parts}
  \item Avoiding preparation
  \item Deep breathing and positive visualization
  \item Avoiding eye contact with the audience
  \item Reading directly from notes without looking up
\end{parts}

\medskip


\noindent   \textbf{Question 70.} Feedback received after a speech practice session is primarily intended to:\pts{1}

\begin{parts}
  \item Discourage the speaker
  \item Identify strengths and areas for improvement
  \item Rate the speaker on a strict numerical scale
  \item Change the topic of the speech
\end{parts}

\medskip

\vfill
\rule{\linewidth}{0.4pt}

\begin{center}\small
\href{https://bhu-exam.vercel.app/}{Click here to visit our website}\\[4pt]
\href{https://me-aryan.vercel.app/}{Click here to visit our contributor}
\end{center}

\end{document}
